% ===========================================================================
% objetivos.tex  --  Capitulo 2
% Se incluye desde main.tex con \input{capitulos/objetivos}
% ===========================================================================
\section{Objetivos}
\label{cap:objetivos}

\subsection{Objetivo general}
\label{sec:obj-general}

Desarrollar un \textit{pipeline} computacional completamente \textit{in silico},
basado en redes neuronales de grafos, que permita \textbf{priorizar} moléculas
candidatas dirigidas al receptor 5-HT\textsubscript{1A} mediante la integración
de tres modelos predictivos —afinidad de unión, permeabilidad a la barrera
hematoencefálica y perfil de toxicidad—, evaluados bajo un protocolo de
validación que refleje el uso previsto del sistema: predecir sobre quimiotipos
no vistos durante el entrenamiento.

\subsection{Objetivos específicos}
\label{sec:obj-especificos}

Cada objetivo se acompaña de un criterio de verificación. La etiqueta entre
paréntesis indica la fase del trabajo terminal en la que se desarrolla.

\begin{enumerate}

  \item \textbf{Construir y curar un conjunto de datos de ligandos del receptor
        5-HT\textsubscript{1A}} a partir de repositorios públicos (ChEMBL,
        BindingDB y PDSP Ki Database), aplicando una política explícita de
        estandarización química —desalinización, neutralización, selección de
        tautómero canónico, tratamiento de mezclas y decisión documentada sobre
        estereoquímica— y una política explícita de deduplicación por InChIKey
        y de agregación de mediciones discordantes. \hfill \textit{(TT1)}

        \textit{Verificación:} tabla de trazabilidad que reporte, por cada
        fuente, la consulta exacta, la fecha de descarga, el número de registros
        iniciales, los eliminados en cada etapa y el número final de compuestos
        únicos.

  \item \textbf{Construir los conjuntos de toxicidad y de permeabilidad de la
        barrera hematoencefálica} (Tox21 y B3DB/BBBP) bajo el mismo
        procedimiento de estandarización, garantizando la ausencia de
        compuestos compartidos entre particiones. \hfill \textit{(TT1)}

        \textit{Verificación:} análisis exploratorio que documente el balance de
        clases por punto final, la distribución de propiedades fisicoquímicas y
        el solapamiento estructural entre conjuntos.

  \item \textbf{Definir y fijar el protocolo de evaluación antes del modelado},
        adoptando partición por \textit{scaffold} de Bemis--Murcko como criterio
        primario para las tres tareas, con partición aleatoria reportada
        únicamente como cota superior de referencia y etiquetada como tal.
        \hfill \textit{(TT1)}

        \textit{Verificación:} semillas fijas, versiones de bibliotecas
        ancladas, y particiones persistidas en disco de modo que cualquier
        modelo posterior se evalúe sobre exactamente los mismos conjuntos.

  \item \textbf{Establecer líneas base fuertes} mediante huellas digitales
        circulares (ECFP4/Morgan) combinadas con bosques aleatorios y modelos de
        \textit{gradient boosting}, con ajuste de hiperparámetros comparable al
        que recibirán las redes neuronales de grafos. \hfill \textit{(TT1)}

        \textit{Verificación:} desempeño reportado con intervalos de confianza
        sobre la partición por \textit{scaffold}. Si la línea base supera a la
        GNN, ese resultado se reporta como hallazgo legítimo.

  \item \textbf{Implementar versiones iniciales de los tres modelos basados en
        redes neuronales de grafos}: regresión de pKi para afinidad,
        clasificación multitarea para toxicidad y clasificación binaria para
        permeabilidad de la BHE. \hfill \textit{(TT1)}

        \textit{Verificación:} regresión evaluada con RMSE, MAE y
        $R^2$; clasificación evaluada con AUPRC, matriz de confusión y
        ROC-AUC, en ese orden de prioridad, dado el desbalance de clases.

  \item \textbf{Declarar el dominio de aplicabilidad} de cada modelo, en
        términos del rango de la variable respuesta y de la similitud
        estructural respecto al conjunto de entrenamiento, e implementar un
        mecanismo que marque las predicciones fuera de dominio.
        \hfill \textit{(TT1--TT2)}

        \textit{Verificación:} distribución de similitud de Tanimoto entre
        consulta y vecinos más próximos del conjunto de entrenamiento, con
        umbral justificado.

  \item \textbf{Optimizar las arquitecturas de grafos y sus hiperparámetros},
        explorando variantes de agregación y de representación de nodos y
        aristas. \hfill \textit{(TT2)}

  \item \textbf{Integrar los tres modelos en un único \textit{pipeline}} con un
        criterio de priorización explícito que combine las tres predicciones y
        propague la incertidumbre de cada una. \hfill \textit{(TT2)}

  \item \textbf{Aplicar el \textit{pipeline} integrado sobre una biblioteca de
        compuestos virtuales} para proponer un conjunto ordenado de candidatos.
        \hfill \textit{(TT2)}

        \textit{Verificación:} los candidatos propuestos se reportan como
        \emph{hipótesis priorizadas}, acompañadas de su posición respecto al
        dominio de aplicabilidad. \todo{Especificar la biblioteca concreta
        (ZINC, Enamine REAL u otra) y su tamaño.}

\end{enumerate}

\subsection{Delimitación y no-objetivos}
\label{sec:obj-delimitacion}

Declarar lo que el trabajo \textbf{no} hace es tan necesario como declarar lo que
hace, porque acota el criterio con el que debe juzgarse.

\begin{itemize}
  \item No se realiza síntesis química ni validación experimental
        \textit{in vitro} o \textit{in vivo} de ningún compuesto. Las salidas
        del sistema son predicciones, no evidencia de actividad.
  \item No se emplea información estructural tridimensional del receptor ni
        acoplamiento molecular (\textit{docking}). El enfoque es basado en
        ligando.
  \item No se modela selectividad frente a otros subtipos de receptores
        serotoninérgicos ni perfiles de polifarmacología.
  \item No se predice eficacia clínica, dosis ni farmacocinética completa. La
        permeabilidad de la BHE se trata como una clasificación binaria y no
        como un modelo de distribución.
  \item No se afirma relación causal alguna entre las subestructuras
        identificadas y la actividad observada.
\end{itemize}

\begin{quote}\small
\textbf{Delimitación TT1/TT2.} El \textbf{TT1} abarca los objetivos 1 a 6:
construcción y curación de los datos, fijación del protocolo de evaluación,
líneas base y modelado inicial de referencia. El \textbf{TT2} abarca los
objetivos 7 a 9: optimización de arquitecturas, integración del
\textit{pipeline} y cribado de la biblioteca virtual.
\end{quote}